\chapter{诊断}

% === 源文件: diagnostics/flags.md ===
\section{诊断标志}


诊断标志让你可以启用定向调试日志，而无需在所有地方开启详细日志。标志是可选启用的，除非子系统检查它们，否则不会生效。

\subsection{工作原理}


\begin{itemize}
  \item 标志是字符串（不区分大小写）。
  \item 你可以在配置中或通过环境变量覆盖来启用标志。
  \item 支持通配符：
  \item \texttt{telegram.*} 匹配 \texttt{telegram.http}
  \item \texttt{*} 启用所有标志
\end{itemize}


\subsection{通过配置启用}


\begin{lstlisting}[language=json]
{
  "diagnostics": {
    "flags": ["telegram.http"]
  }
}
\end{lstlisting}


多个标志：

\begin{lstlisting}[language=json]
{
  "diagnostics": {
    "flags": ["telegram.http", "gateway.*"]
  }
}
\end{lstlisting}


更改标志后重启 Gateway 网关。

\subsection{环境变量覆盖（一次性）}


\begin{lstlisting}[language=bash]
OPENCLAW_DIAGNOSTICS=telegram.http,telegram.payload
\end{lstlisting}


禁用所有标志：

\begin{lstlisting}[language=bash]
OPENCLAW_DIAGNOSTICS=0
\end{lstlisting}


\subsection{日志存储位置}


标志将日志输出到标准诊断日志文件。默认位置：

\begin{lstlisting}
/tmp/openclaw/openclaw-YYYY-MM-DD.log
\end{lstlisting}


如果你设置了 \texttt{logging.file}，则使用该路径。日志为 JSONL 格式（每行一个 JSON 对象）。脱敏仍然根据 \texttt{logging.redactSensitive} 应用。

\subsection{提取日志}


选择最新的日志文件：

\begin{lstlisting}[language=bash]
ls -t /tmp/openclaw/openclaw-*.log | head -n 1
\end{lstlisting}


过滤 Telegram HTTP 诊断：

\begin{lstlisting}[language=bash]
rg "telegram http error" /tmp/openclaw/openclaw-*.log
\end{lstlisting}


或在复现时使用 tail：

\begin{lstlisting}[language=bash]
tail -f /tmp/openclaw/openclaw-$(date +%F).log | rg "telegram http error"
\end{lstlisting}


对于远程 Gateway 网关，你也可以使用 \texttt{openclaw logs --follow}（参见 \href{/cli/logs}{/cli/logs}）。

\subsection{注意事项}


\begin{itemize}
  \item 如果 \texttt{logging.level} 设置为高于 \texttt{warn}，这些日志可能会被抑制。默认的 \texttt{info} 级别即可。
  \item 标志可以安全地保持启用状态；它们只影响特定子系统的日志量。
  \item 使用 \href{/logging}{/logging} 更改日志目标、级别和脱敏设置。
\end{itemize}



\clearpage
